\section*{الفصل \ref{ch:b3:lab-techniques-3} --- التقنيات المخبرية الثالثة: ممارسة البحث}

\begin{solution}{exo:b3:lab-techniques-3:1}
$(0.50/1013)^3 = \num{1.2e-10}$. فالتسربات، والغاز الممتز على الزجاج، وانبعاث الغازات من الشحم والحواجز المطاطية
تترك أكثر من ذلك بكثير؛ والأدوات الزجاجية الجافة الساخنة والوصلات الجيدة أهم من
الدورات الإضافية.
\end{solution}

\begin{solution}{exo:b3:lab-techniques-3:2}
$\qty{20.0}{mmol}/\qty{1.48}{mol/L} = \qty{13.5}{mL}$.
\end{solution}

\begin{solution}{exo:b3:lab-techniques-3:3}
$120 + 10 \times 1.00782503 + 55.93493633 - 0.00054858 = \qty{186.01264}{u}$. الخطأ: $10^6 \times (186.0129 - 186.01264)/186.01264 = +1.4$ ppm،
وهو ضمن حد 5 ppm بفارق مريح.
\end{solution}

\begin{solution}{exo:b3:lab-techniques-3:4}
الأولية: المقالة في المجلة، وبراءة الاختراع، والأطروحة. الثانوية: المقالة الاستعراضية، 
والكتاب المرجعي للثوابت. الثالثية: الكتاب الدراسي.
\end{solution}

\begin{solution}{exo:b3:lab-techniques-3:5}
$M = \qty{194.0}{g/mol}$: C $96.0/194.0 = 49.48$\,\%، وH $10.0/194.0 = 5.15$\,\%، وN $56.0/194.0 = 28.87$\,\%. والفروق 0.17 و0.07 و0.16
نقطة: ضمن 0.4، فالتحليل يدعم الصيغة.
\end{solution}

\begin{solution}{exo:b3:lab-techniques-3:6}
$\qty{170.0}{mg}/\qty{212.0}{g/mol} = \qty{0.802}{mmol}$؛ $c = \qty{0.802}{mmol}/\qty{0.54}{mL} = \qty{1.5}{mol/L}$ (برقمين معنويين، يفرضهما
الحجم المقروء).
\end{solution}

\begin{solution}{exo:b3:lab-techniques-3:7}
المتوسطان 73.0 و77.0\,\%؛ والانحرافان المعياريان كلاهما $\sqrt{10/4} = 1.58$؛ و$s_p = 1.58$. و$t = 4.0/(1.58\sqrt{2/5}) = 4.0$، أكبر من
2.31: فالتحسن دال عند المستوى 95\,\%.
\end{solution}

\begin{solution}{exo:b3:lab-techniques-3:8}
$F = (0.80/0.30)^2 = 7.1 > 5.05$: الطريقة الأولى أقل دقة على نحو دال.
\end{solution}

\begin{solution}{exo:b3:lab-techniques-3:9}
اشترِ قوارير صغيرة فيها مثبّط، وأرّخ كلًا منها عند الاستلام وعند الفتح، وخزّنها
مغلقة في مكان مظلم بارد؛ واختبرها بحثًا عن البيروكسيدات قبل أي تقطير أو تبخير
وعلى فترات ثابتة بعد الفتح؛ وتخلّص من القوارير المنتهية التاريخ، أو
ذات الاختبار الإيجابي، عبر مسار النفايات الخطرة دون تركيزها.
\end{solution}

\begin{solution}{exo:b3:lab-techniques-3:10}
الأخطار: انطلاق الهيدروجين (حريق، وضغط)، وتفاعلية NaH مع الماء،
وتفكك ناشر للحرارة معروف للمذيب DMF مع هيدريد الصوديوم قد ينفلت
عند التسخين؛ وDMF أيضًا سام للتكاثر. الضوابط: استبدل المذيب
(THF أو 2-ميثيل رباعي هيدروفوران) أو القاعدة؛ وإن أُبقي عليه، فنفّذ التجربة على نطاق صغير أولًا، عند درجة حرارة
منخفضة، مضيفًا NaH على دفعات مع متابعة درجة الحرارة، تحت ساحبة الأبخرة
خلف درع، مع التنفيس عبر فقّاعة، وحمام تبريد جاهز؛ وإجراء
مكتوب، وحضور شخص ثانٍ، ومعدات الوقاية في آخر القائمة.
\end{solution}

\begin{solution}{exo:b3:lab-techniques-3:11}
المساهمات النسبية: 0.50، و0.50، و$0.010/20.0 = 0.05$، و$0.010/15.0 = 0.067$ و0.10\,\%؛ ومجموعها
$\sqrt{0.25 + 0.25 + 0.0025 + 0.0044 + 0.010} = 0.72$\,\%. والتكاملان هما الغالبان: فنسبة إشارة إلى ضجيج أفضل
وتصحيح دقيق لخط الأساس والطور أهم من الميزان.
\end{solution}

\begin{solution}{exo:b3:lab-techniques-3:12}
$Q = (89.0 - 82.0)/(89.0 - 80.6) = 0.83 > 0.71$: الاختبار ينبّه إليها. وقبل رفضها، ابحث في \omterm{def:b3:lab-techniques-3:notebook}{دفتر المختبر}
عن سبب (قراءة خاطئة للسحاحة، أو فقاعة هواء، أو دفعة مختلفة من محلول المعايرة)؛ فإن
وُجد، فارفضها واذكر السبب؛ وإلا فاحتفظ بها، وأبلغ عن الاختبار، أو أعد
القياس.
\end{solution}

\begin{solution}{pb:b3:lab-techniques-3:1}
\textbf{1.} $\ce{FeCl2 + 2NaC5H5 -> Fe(C5H5)2 + 2NaCl}$.

\textbf{2.} يُتلَف سيكلوبنتاديينيد الصوديوم بالأكسجين والماء، ويتأكسد كلوريد الحديد(II)
إلى الحديد(III) في الهواء الرطب.

\textbf{3.} $\qty{5.00}{g}/\qty{126.8}{g/mol} = \qty{0.0394}{mol}$؛ $\times \qty{185.8}{g/mol} = \qty{7.33}{g}$؛ والمردود $5.86/7.33 = 80$\,\%.

\textbf{4.} $(0.10/1000)^3 = 10^{-12}$ (نظريًا).

\textbf{5.} بتسخين الثنائي، الذي يخضع لتفاعل ديلز--ألدر عكسي، 
وتقطير المونومر المتطاير؛ وعند درجة حرارة الغرفة يعود المونومر إلى التثنّي
بتفاعل ديلز--ألدر، فيُحفظ باردًا ويُستعمل فورًا.

\textbf{6.} كل دفعة تطلق هيدروجينًا وحرارة: فإضافتها ببطء تُبقي تدفق
الغاز ودرجة الحرارة تحت السيطرة.

\textbf{7.} \qty{186.0126}{u}؛ والخطأ $+1.4$ ppm.

\textbf{8.} C $120.0/185.8 = 64.58$\,\%، وH $10.0/185.8 = 5.38$\,\%؛ والمقيس 64.41 و5.47: الفرقان $-0.17$ و$+0.09$، ضمن 0.4.

\textbf{9.} ذرات الهيدروجين العشر متكافئة بالتناظر، والحلقتان تدوران بسرعة.

\textbf{10.} إشارة واحدة، لأن ذرات الكربون العشر كلها متكافئة.

\textbf{11.} بنية بالأشعة السينية لبلورة مفردة.

\textbf{12.} لا: فهو، بوصفه معقدًا ذا 18 إلكترونًا (\cref{ch:b3:organometallic-bonding})، مستقر في الهواء، خلافًا
لسلائفه.

\textbf{13.} لا تتداخل إشاراته مع إشارات الفيروسين، وهو صلب نقي، غير متطاير،
وغير مسترطب يمكن وزنه بدقة، وخامل تجاه
العينة.

\textbf{14.} استرخاء كامل لكل إشارة بين المسوح (تأخير يساوي عدة أضعاف $T_1$)،
وإثارة منتظمة، ونسبة إشارة إلى ضجيج كافية، وتصحيح دقيق للطور وخط
الأساس.

\textbf{15.} الفيروسين $\ce{C10H10Fe}$ \qty{185.8}{g/mol}؛ و1،3،5-ثلاثي ميثوكسي بنزين $\ce{C9H12O3}$ \qty{168.0}{g/mol}.

\textbf{16.} $P_x = 3.942 \times (3/10) \times (185.8/168.0) \times (15.0/20.0) \times 0.999 = 0.980$: أي 98.0\,\%.

\textbf{17.} النسبي $\sqrt{0.50^2 + 0.50^2 + 0.05^2 + 0.067^2 + 0.10^2} = 0.72$\,\%؛ والمطلق $0.72\,\% \times 98.0 = 0.7$ نقطة مئوية.

\textbf{18.} نعم، لكن على نحو ضعيف: فشائبة بنسبة 2\,\% ذات تركيب مشابه تغيّر C وH بأقل
مما يستطيع التحليل كشفه؛ \omterm{def:b3:lab-techniques-3:qnmr}{والرنين المغناطيسي النووي الكمي} هو القياس الأغنى معلومات عن النقاوة.

\textbf{19.} هيدريد الصوديوم (هيدروجين مع الماء، وحريق)، والهيدروجين المنطلق، وTHF 
والسيكلوبنتاديين (شديدا الاشتعال؛ وTHF \omterm{def:b3:lab-techniques-3:peroxide}{مكوِّن للبيروكسيدات})، والتكسير الساخن
للثنائي، وكلوريد الحديد(II) (مهيّج)، والفيروسين (صلب قابل للاشتعال، ضار).

\textbf{20.} هيدريد الصوديوم: يوزن في صورة معلَّق في الزيت تحت غاز خامل، ويُضاف على
دفعات إلى محلول مبرَّد مقلَّب، ويُنفَّس الغاز عبر فقّاعة، ولا ماء
قريبًا. وTHF: جاف وخالٍ من البيروكسيدات من عمود، ويُستعمل تحت ساحبة الأبخرة تحت النيتروجين.

\textbf{21.} تحت غاز خامل ومع التبريد، بإضافة بطيئة للإيزوبروبانول، ثم
الإيثانول، ثم الماء، مع الانتظار في كل مرة حتى يتوقف انطلاق الغاز.

\textbf{22.} موجة عكوسة بإلكترون واحد، $\ce{Fe(C5H5)2+}/\ce{Fe(C5H5)2}$، بقمتين يفصل بينهما نحو 59 mV عند \qty{25}{\celsius}
(\cref{ch:b3:electrode-kinetics})؛ والمزدوجة سريعة ومستقرة وتكاد لا تتأثر
بالمذيب، وهذا يجعلها مرجعًا داخليًا ملائمًا.

\textbf{23.} \omterm{def:b3:lab-techniques-3:notebook}{دفتر المختبر}: التاريخ، والمخطط، والكميات ومصادرها، \omterm{def:b3:lab-techniques-3:assessment}{وتقييم المخاطر}، وما
أُنجز وما لوحظ، والمردودات وأسماء ملفات \omterm{def:b3:lab-techniques-3:notebook}{البيانات الخام}. \omterm{def:b3:lab-techniques-3:literature}{والمعلومات
الداعمة}: الإجراء الكامل وكل بيانات التوصيف مع نسخ من
الأطياف.

\textbf{24.} كي يتمكن الآخرون، والمؤلفون أنفسهم بعد سنوات، من فتحها ومعالجتها من جديد
أيًا كانت البرمجيات المتوفرة لديهم: قابلة للإيجاد، وقابلة للوصول، وقابلة للتشغيل البيني، وقابلة لإعادة الاستعمال.

\textbf{25.} نقاوة عينة الفيروسين \omterm{def:b3:lab-techniques-3:qnmr}{بالرنين المغناطيسي النووي الكمي} هي $98.0 \pm 0.7$\,\% (ارتياب معياري).
\end{solution}
